\documentclass[11pt,letterpaper,openany]{book}
\usepackage{beautybook-inspired}

\newcommand{\BookTitle}{Building a CSV Profiling and ESP32 OTA Dashboard}
\newcommand{\BookSubtitle}{A friendly, hands-on guide from first click to first update}
\newcommand{\BookAuthor}{Adji-SP}
\hypersetup{
  pdftitle={Building a CSV Profiling and ESP32 OTA Dashboard},
  pdfauthor={Adji-SP},
  pdfsubject={Beginner guide to the CSV Profiling dashboard and ESP32-S3 OTA system}
}

\begin{document}
\frontmatter

\bookcover{\BookAuthor}{Building a CSV Profiling\\and ESP32 OTA Dashboard}{\BookSubtitle}{First Edition}{CSV\_PROFILING}

\cleardoublepage
\thispagestyle{empty}
\vspace*{0.58\textheight}
{\small
\textbf{\BookTitle}\par
First edition, September 2026.\par

This handbook documents the local project as implemented in this repository. The supplied \emph{beautybook} PDF inspired its visual system. Its text and examples were written specifically for CSV\_PROFILING.\par

Commands that affect hardware are marked clearly. Review the ESP32-S3 flash size, serial port, partition table, and local network settings before using them.\par
}
\clearpage

\frontbanner{Preface}
Open the project folder and it can feel as though three different subjects have landed on the same desk: data analysis, web development, and embedded Rust. That feeling is normal. We will pick up one piece at a time.

The dashboard has two working rooms behind it. Flask and Pandas look after CSV files. Axum looks after firmware builds and devices. Farther down the line, an ESP32-S3 calls home over Wi-Fi and asks whether a new application is waiting. Once you can picture those rooms, the source tree stops looking like a wall of unfamiliar names.

This book is meant to be used with the project open beside it. Read a few pages, touch the file being discussed, and try the small checkpoint at the end of the chapter. Keep the terminals where you can see them. A failed command often tells you exactly which door is still locked.

\begin{flushright}
\textit{The CSV\_PROFILING project}\par
2026
\end{flushright}
\clearpage

\frontbanner{How to Use This Book}
Treat these pages like a marked-up workshop notebook. You do not need to memorize the code. Keep the repository open, follow the paths printed in monospace, and compare the short examples with the real files. Windows commands use PowerShell. Linux and macOS commands use a POSIX shell.

You will meet three kinds of notes in the margins of the lesson:

\begin{mentorbox}
A tutor's note slows down for the question someone would usually ask across the table.
\end{mentorbox}

\begin{warningbox}
A pause note appears before a command can touch hardware, delete data, or loosen a safety boundary.
\end{warningbox}

\begin{checkpoint}
A small exercise lets your own screen confirm what the page just explained.
\end{checkpoint}

You can begin without an ESP32 on the desk. The dashboard, APIs, database, ZIP checks, and error messages all run on the computer. When the lesson reaches a physical flash, the page will say so plainly.
\clearpage

\tableofcontents

\mainmatter

\part{Foundations}
\partintro{We begin with the view from the doorway. You will see the whole project in plain language, then start it without rushing past the small signs that tell us it is healthy.}
\chapter{Meet the Project Before We Touch It}

\section{One dashboard, two helpers}

Open the dashboard in your browser. It looks like one application because every page shares the same sidebar, colors, and buttons. Behind that screen sit two separate helpers.

The Python helper listens on port 5000. Give it a CSV file and it will inspect the rows, calculate summaries, build a report, and apply cleaning choices. The Rust helper listens on port 7000. Give it a firmware project and it will guard the upload, run the build tools, save the finished binary, and remember which device should receive it.

Picture a small repair shop with one front desk and two workbenches. The front desk is the browser. One bench handles data. The other handles firmware.

\begin{center}
\begin{tcolorbox}[width=0.9\textwidth,colback=BookPaper,colframe=BookTealLight,arc=3pt]
\centering\ttfamily
You, using the dashboard at :5000\par
\vspace{0.25em}
$\swarrow$\hspace{7em}$\searrow$\par
\vspace{0.25em}
Flask + Pandas :5000\hspace{3em}Axum :7000\par
CSV work\hspace{8em}firmware work\par
\hspace*{13em}$\downarrow$\par
\hspace*{11em}ESP32-S3 on Wi-Fi
\end{tcolorbox}
\end{center}

Keeping the benches separate protects the CSV profiler. Adding a firmware feature does not require us to rebuild the data service in Rust. The browser only needs to know which address to call.

\begin{mentorbox}[A name you will see often]
Developers call the handoff between two programs a \emph{service boundary}. Here the boundary is simply an HTTP request. The phrase sounds formal; the idea is ordinary. One program asks another program to do a clearly named job.
\end{mentorbox}

\section{Follow a CSV file with your finger}

Suppose you choose \codepath{readings.csv} and press Run Profiling. The browser sends the file to Flask. Flask saves a working copy and starts the slower profiling work in the background. While Pandas reads the data, the browser receives progress messages. When the report is ready, Flask serves the finished HTML page.

Every later CSV action returns to the same helper. Previous Runs asks Flask for saved sessions. Cleaning asks Flask to change a working copy. The playground asks Flask to perform a small Pandas operation.

You can remember the route in one line:

\begin{center}
\ttfamily browser $\rightarrow$ Flask $\rightarrow$ Pandas $\rightarrow$ report
\end{center}

\section{Now follow a firmware ZIP}

The OTA journey begins with a ZIP because a Rust program is more than one source file. It may have a Cargo manifest, extra modules, and library dependencies. The browser sends that ZIP, a board choice, and a version to Axum.

Axum checks the archive before opening it. Then it asks Cargo to compile the project for the ESP32-S3. Cargo produces an ELF executable. \codepath{espflash} turns that ELF into an ESP-IDF application image. The server calculates a SHA-256 fingerprint and saves the binary under \codepath{firmware-builds}.

Later, the device asks whether an update is waiting. If one is assigned, the ESP32 downloads the \codepath{.bin}. It never receives the ZIP or a loose \codepath{.rs} file.

\begin{center}
\ttfamily ZIP $\rightarrow$ Cargo $\rightarrow$ ELF $\rightarrow$ application .bin $\rightarrow$ ESP32
\end{center}

\section{Use the folders as signposts}

Keep this table nearby during your first few sessions. You only need the first sentence for each folder.

\begin{longtable}{>{\ttfamily\footnotesize\raggedright\arraybackslash}p{0.39\textwidth}p{0.53\textwidth}}
\toprule
Folder or file & What lives there\\
\midrule
\endhead
web/ & The pages you see: HTML, CSS, and browser JavaScript.\\
python/profiler\_server.py & Flask routes and all existing CSV work.\\
ota-server/ & The Axum server, database code, ZIP checks, build runner, and device API.\\
firmware-projects/ & Private working folders made when Axum extracts uploads.\\
firmware-builds/ & Finished OTA application binaries.\\
data/ota.db & The SQLite notebook where Axum remembers projects, builds, firmware, devices, and deployments.\\
examples/esp32-rust-ota-client/ & The steady base program installed on the board and reused for managed builds.\\
examples/uploaded-firmware-basic/ & The smallest application ZIP you can study and upload.\\
examples/uploaded-firmware-with-libs/ & A second application showing extra Cargo libraries.\\
docs/OTA\_GUIDE.md & Shorter engineering notes about partitions, rollback, and networking.\\
\bottomrule
\end{longtable}

The three storage folders fill themselves while the system runs. Your hand-written source belongs in \codepath{web}, \codepath{python}, \codepath{ota-server}, or an application project under \codepath{examples}.

\section{Which main program is the real main program?}

There are three programs running in three places. The browser runs JavaScript. The PC runs Flask and Axum. The ESP32 runs one compiled firmware image.

Inside that firmware image, \codepath{examples/esp32-rust-ota-client/src/main.rs} owns the true executable \codepath{fn main()}. Your uploaded application also keeps a file named \codepath{src/main.rs}, but it exports two callable functions named \codepath{setup} and \codepath{main}. During the build, Rust links those functions into the stable runtime.

Think of the runtime as a house with its wiring already installed. Your uploaded application decides what happens in the rooms. The final build produces one complete house. No second program balances on top of it.

\section{Why the first flash uses USB}

A new ESP32 does not know your Wi-Fi name. It has no address for the OTA server and no two-slot OTA partition table. USB gives it that first foundation: bootloader, partitions, Wi-Fi-aware runtime, and a small default application.

After the first successful flash, the board can introduce itself to the server and ask for later releases over Wi-Fi.

\begin{mentorbox}[The short answer]
Yes, flash the stable runtime to the ESP32-S3 first. You normally do this once while provisioning the board. OTA can take over after the device can join the network and reach port 7000 on your PC.
\end{mentorbox}

\section{Who gets the Wi-Fi modem?}

The ESP32 has a single modem resource. The runtime takes it once and gives it to \codepath{wifi-manager}. That manager keeps the connection alive for heartbeats and OTA checks. Your application receives the network's current state through \codepath{AppContext}, along with the remaining pins and hardware controllers it may use.

This arrangement leaves room for a friendlier Wi-Fi setup screen later. Saved credentials, reconnect rules, or a captive portal can stay inside the manager. Sensor applications can remain focused on sensors.

\begin{warningbox}
Do not call \codepath{Peripherals::take()} inside an uploaded application. The runtime has already taken the global set. Reach for \codepath{context.peripherals} and take only the pin or controller your application needs.
\end{warningbox}

\begin{checkpoint}
Open the repository and find these four places: \codepath{web/index.html}, \codepath{python/profiler_server.py}, \codepath{ota-server/src/main.rs}, and \codepath{examples/esp32-rust-ota-client/src/main.rs}. Say aloud which machine runs each one. That small sentence prevents a great deal of confusion later.
\end{checkpoint}

\chapter{Let's Get It Running}

\section{Begin with the computer}

Leave the ESP32 unplugged for the moment. The first win is smaller: open the dashboard and make both PC services answer. Once that works, every later problem has fewer hiding places.

You need Python 3.10 or newer for the CSV side and Rust for the OTA side. The embedded toolchain can wait until the dashboard is comfortable.

\section{Prepare Python gently}

A virtual environment gives this project its own box of Python packages. It keeps your other Python work out of the way.

Open PowerShell in the repository and type the first line:

\begin{lstlisting}
py -3 -m venv python\.venv
\end{lstlisting}

When the command returns, look inside \codepath{python}. A new \codepath{.venv} folder should be there. Now install the packages named in \codepath{python/requirements.txt}:

\begin{lstlisting}
python\.venv\Scripts\python -m pip install -r python\requirements.txt
\end{lstlisting}

On Linux or macOS, the same two steps are:

\begin{lstlisting}
python3 -m venv python/.venv
python/.venv/bin/pip install -r python/requirements.txt
\end{lstlisting}

The installation may take a while because data-profiling libraries bring several helpers with them. Wait for the command prompt to return. Red text near the middle can be a warning; the final lines tell you whether installation finished or failed.

\section{Check the Rust tools you already have}

Run one command at a time:

\begin{lstlisting}
cargo --version
rustc --version
\end{lstlisting}

If both print a version, the PC can begin building the Axum server. Firmware needs two more checks later:

\begin{lstlisting}
rustc +esp --version
espflash --version
\end{lstlisting}

Do not chase every missing embedded tool yet. The dashboard can still start, store projects, validate ZIP files, and show an honest failed-build message.

\begin{mentorbox}[A Windows detail that surprises many people]
The ESP Rust toolchain may be installed while Windows still lacks \codepath{link.exe}. That file comes with Visual Studio Build Tools and the Desktop development with C++ workload. If a build names \codepath{link.exe}, the Rust code is not necessarily at fault. The compiler is asking Windows for a native tool it cannot find.
\end{mentorbox}

\section{Make your private .env file}

Copy \codepath{.env.example} and name the copy \codepath{.env}. This is your local settings page. It may hold a Wi-Fi password, so Git should ignore it.

For the first start, pay attention to these lines:

\begin{lstlisting}[language=TOML,numbers=none]
OTA_BIND_ADDRESS=0.0.0.0
OTA_PORT=7000
OTA_PUBLIC_BASE_URL=http://192.168.1.5:7000

WIFI_SSID=your-wifi-name
WIFI_PASS=your-wifi-password
DEVICE_NAME=ESP32-S3 Managed Device
\end{lstlisting}

Replace \codepath{192.168.1.5} with the LAN address of your own PC. On Windows, run \codepath{ipconfig} and look under the Wi-Fi or Ethernet adapter you are actually using. The useful address often begins with \codepath{192.168} or \codepath{10}.

\begin{warningbox}[Why localhost will betray the device]
On your PC, \codepath{localhost} means your PC. On the ESP32, the same word means the ESP32. Put the PC's LAN address in \codepath{OTA_PUBLIC_BASE_URL} and in the runtime's \codepath{OTA_SERVER_URL}.
\end{warningbox}

The rest of \codepath{.env.example} already points to sensible local folders. Leave those values alone during your first run.

\section{Press the project's start button}

Windows users can run:

\begin{lstlisting}
.\run-dev.bat
\end{lstlisting}

Two terminals should open. One says CSV Profiler and uses port 5000. The other compiles and starts the Rust OTA server on port 7000. The first Rust build may be quiet for a while because Cargo is preparing dependencies.

Linux and macOS users can run:

\begin{lstlisting}
chmod +x run-dev.sh
./run-dev.sh
\end{lstlisting}

When both helpers are ready, open \codepath{http://localhost:5000}. You should see the sidebar and the CSV upload page.

\section{Ask each service a tiny question}

The dashboard itself proves Flask is answering. Give Axum a similarly small test by opening:

\begin{lstlisting}[numbers=none]
http://localhost:7000/api/ota/status
\end{lstlisting}

The browser should show JSON. Do not worry about every field yet. Find \codepath{"status":"online"} and check that \codepath{public_base_url} carries your LAN address.

\begin{lstlisting}[language=JSON,numbers=none]
{
  "service": "CSV_PROFILING OTA Server",
  "status": "online",
  "public_base_url": "http://192.168.1.5:7000",
  "supported_targets": ["esp32s3"]
}
\end{lstlisting}

\section{If the page refuses to open}

Do not restart everything at once. Look at one door at a time.

First, visit \codepath{http://127.0.0.1:5000}. If it also fails, read the Flask terminal. A missing Python package usually names itself there.

If the dashboard opens but its OTA light says offline, visit the status URL on port 7000. Read the Rust terminal if that address fails. Cargo may still be compiling, another program may own the port, or configuration may have stopped startup.

If both addresses work on the PC and the ESP32 later cannot connect, move your attention to the LAN address, firewall, and Wi-Fi network. That is a different door.

Open the dashboard through HTTP. Double-clicking \codepath{web/index.html} creates a \codepath{file:///} page with a different browser origin and can make good API code look broken.

\begin{checkpoint}
Keep two tabs open. The first shows \codepath{http://localhost:5000}. The second shows Axum's status JSON. Point to the two matching terminal windows. You now have a clean starting line, even if no ESP32 is connected.
\end{checkpoint}


\part{The Dashboard and CSV Profiler}
\partintro{A single CSV file will be our guide. We will follow it from the upload button into Python, watch the report come back, and peek behind each dashboard screen as we go.}
\chapter{Walk Through the Dashboard}

\section{Start with what your eyes can see}

Keep the dashboard open beside the project folder. The dark strip on the left is the sidebar. The thin bar across the top tells you which page is active. The large pale area is where the current task appears.

Click CSV Profiler, Previous Runs, OTA Firmware, Devices, and Build History. The browser does not load five separate HTML pages. It already has each section and shows one at a time.

Now open \codepath{web/index.html}. Search for \codepath{ota-section}. You will find the OTA Firmware screen. A little farther down are \codepath{devices-section} and \codepath{builds-section}. The HTML is the dashboard's floor plan.

\begin{mentorbox}[A useful beginner habit]
When a button behaves strangely, find its \codepath{id} in the HTML first. Then search that same ID in JavaScript. This is often faster than reading an entire script from the top.
\end{mentorbox}

\section{Meet the frontend files one by one}

Only four files need our attention:

\begin{description}
  \item[\codepath{web/index.html}] The objects on the page: links, headings, forms, tables, dialogs, and empty states.
  \item[\codepath{web/style.css}] Their appearance: spacing, color, borders, cards, buttons, badges, progress bars, and the dark build terminal.
  \item[\codepath{web/app.js}] The behavior of CSV pages.
  \item[\codepath{web/ota.js}] The behavior of OTA Firmware, Devices, and Build History.
\end{description}

The HTML loads both scripts near its closing tag:

\begin{lstlisting}[language=HTML,numbers=none]
<script src="app.js"></script>
<script src="ota.js"></script>
\end{lstlisting}

That order lets the old CSV script keep its place. OTA behavior arrives beside it in a separate file.

\section{A button is a short chain of events}

Look at the Upload \& Build button in \codepath{index.html}. Its ID is \codepath{ota-build-btn}. Search for that text in \codepath{ota.js}. You will find the event handler that reads the chosen file, target, and version.

The chain looks like this:

\begin{center}
\ttfamily click $\rightarrow$ read form $\rightarrow$ call Axum $\rightarrow$ update the page
\end{center}

That same pattern appears throughout the dashboard. A navigation link changes visible sections. A refresh button calls a list endpoint and rebuilds a table. A deploy button opens a dialog, then sends the selected firmware ID.

You do not have to understand every line before making sense of the flow. Find the HTML object, find its listener, then follow the function it calls.

\section{Why there are two API addresses}

At the top of the scripts you will see:

\begin{lstlisting}[language=JavaScript]
const API_BASE = "http://localhost:5000";
const OTA_API_BASE = "http://localhost:7000";
\end{lstlisting}

These two lines are tiny signposts. CSV requests travel to Flask. Firmware and device requests travel to Axum. Do not merge the constants just because both services run on the same PC.

The OTA script wraps its repeated request work in \codepath{otaRequest}. In plain language, that helper says: send the request, check whether it succeeded, read a useful server error when it failed, and decode JSON when a body is present.

\begin{lstlisting}[language=JavaScript]
async function getDevices() {
  return otaRequest("/api/ota/devices");
}

async function deployFirmware(deviceId, firmwareId) {
  return otaRequest(`/api/ota/devices/${deviceId}/deploy`, {
    method: "POST",
    headers: { "Content-Type": "application/json" },
    body: JSON.stringify({ firmware_id: firmwareId })
  });
}
\end{lstlisting}

The first function is a question. The second is an instruction. Both travel through the same careful messenger.

\section{Watch the OTA build page change mood}

Choose a ZIP and press Upload \& Build. Even before the compiler starts, the page moves through small states:

\begin{description}
  \item[Idle] Nothing has been submitted.
  \item[Uploading] The browser is sending the archive.
  \item[Queued] Axum has made a build record.
  \item[Building] The tools on the PC are working.
  \item[Success] A real binary and fingerprint were saved.
  \item[Failed] Something stopped the build, and the log tells us where.
\end{description}

The frontend asks for the build record every few seconds. This is called polling. Imagine looking through the workshop window now and then instead of keeping someone on the phone for the whole build.

Compiler errors remain in the terminal panel. A toast may say the build failed, but the dark log area carries the useful details. Scroll upward to find the first clear error.

\section{Understand the device cards}

The Devices page is a view of what Axum remembers. It does not scan the room for Wi-Fi boards. A real ESP32 puts itself on that page by registering and sending heartbeats.

Each card shows its last reported version, desired version, network address, last-seen time, and OTA progress. If a board loses power, its old card may remain. The timestamp tells you how fresh the information is.

\section{Make one harmless visual change}

Open \codepath{web/index.html} and find the subtitle under OTA Firmware. Change a few words, save the file, and refresh the dashboard. If the old words remain, use a hard refresh so the browser stops using its cached copy.

Then undo your practice wording. You have just traced the simplest frontend loop: edit HTML, let Flask serve it, and let the browser draw it.

\begin{mentorbox}[About the word terminal]
The dark OTA terminal shows build output from the PC. It does not yet show messages printed by the ESP32. For now, use \codepath{espflash monitor} for board output. Chapter 12 describes two honest ways to bring those logs into the dashboard later.
\end{mentorbox}

\begin{checkpoint}
Pick the Devices navigation link. Find its ID in \codepath{index.html}, its click listener in JavaScript, the function that loads devices, and the Axum route it calls. Write those four names on paper as a small trail.
\end{checkpoint}


\chapter{Follow One CSV File}

\section{Make a tiny file we can understand by sight}

Create \codepath{sensor-readings.csv} with these four rows:

\begin{lstlisting}[numbers=none]
device,temperature,room
alpha,24.7,lab
beta,,office
alpha,24.7,lab
gamma,31.2,workshop
\end{lstlisting}

We already know its secrets. One temperature is missing. The alpha row is repeated. The workshop is warmer than the other rooms. A tiny file lets us check whether the report is telling the truth.

Open CSV Profiler, choose the file, and press Run Profiling. Keep \codepath{python/profiler_server.py} open while the progress screen moves.

\section{The upload route receives the envelope}

Search the Python file for \codepath{@app.route("/upload")}. Flask calls the function below that line when the browser posts a file.

Read the first checks as ordinary questions:

\begin{lstlisting}[language=Python]
if "file" not in request.files:
    return jsonify({"error": "No file provided"}), 400

if not f.filename.lower().endswith(".csv"):
    return jsonify({"error": "Only CSV files are supported"}), 400
\end{lstlisting}

If the envelope contains a CSV, Flask gives the job a random ID and saves a working copy. Your original filename goes into the session notes; the stored filename uses the random ID so two people can upload \codepath{data.csv} without choosing the same path.

The route then starts \codepath{run_profiling} on a background thread and answers the browser with the job ID.

\begin{mentorbox}[Why answer before the report is finished?]
A full profile can take time. If one browser request waits silently for the whole job, the page feels frozen. Returning the job ID gives the browser something to follow while Python continues its work.
\end{mentorbox}

\section{Progress arrives like notes under a door}

Every running job has a small Python queue. The profiler places progress messages into it. The route \codepath{/progress/<job-id>} waits for those messages and sends them to the browser as they arrive.

This one-way stream is called server-sent events, usually shortened to SSE. The name matters less than the picture: Python writes short notes; the browser reads them without asking the same question again and again.

\begin{lstlisting}[language=Python]
def generate():
    while True:
        message = q.get()
        if message is None:
            break
        yield message
\end{lstlisting}

When the queue receives \codepath{None}, the stream closes. The browser then opens the finished report.

\section{Where the report sleeps}

Flask stores the report HTML in \codepath{python/reports}. Beside it sits a JSON sidecar. Think of the sidecar as the label on a folder: original filename, job ID, creation time, report address, row count, and column count.

Previous Runs reads those labels through \codepath{GET /sessions}. A browser refresh does not erase the session because the important details live on disk, not only in JavaScript memory.

These routes form the small session shelf:

\begin{longtable}{p{0.15\textwidth}p{0.34\textwidth}p{0.41\textwidth}}
\toprule
Method & Route & What you receive\\
\midrule
\endhead
GET & \codepath{/sessions} & All report sessions Flask can still find.\\
GET & \codepath{/session/<job-id>} & The notes for one saved session.\\
GET & \codepath{/report/<job-id>} & The generated HTML report.\\
DELETE & \codepath{/reset/<job-id>} & Removal of one session and its working files.\\
DELETE & \codepath{/reset-all} & Removal of all local CSV sessions.\\
\bottomrule
\end{longtable}

\section{Clean a copy, then look again}

Return to our four-row file. Ask the cleaning panel to remove duplicates and fill the missing temperature with the median. When you apply the choices, the browser posts a JSON description to \codepath{/clean/<job-id>}.

The idea resembles this:

\begin{lstlisting}[language=JSON,numbers=none]
{
  "drop_duplicates": true,
  "missing": {
    "temperature": "fill_median"
  }
}
\end{lstlisting}

Pandas opens the current working data, applies those choices, and writes \codepath{<job-id>_cleaned.csv}. The original upload remains available. Flask refreshes the report notes and gives the browser a download address.

Open the cleaned CSV. You should see three data rows and no empty temperature. Since \codepath{24.7} appears twice in the original values, the median fill should also be \codepath{24.7}.

\section{Use the playground as a scratch pad}

The playground runs one small Pandas idea at a time. Start with Preview, Head, Tail, Describe, or Value Counts. Those choices do not change the saved data.

For example, this request asks for the ten rows with the largest temperature values:

\begin{lstlisting}[language=JSON,numbers=none]
{
  "operation": "sort_values",
  "column": "temperature",
  "limit": 10,
  "mutate": false
}
\end{lstlisting}

The word \codepath{mutate} means ``change the working data.'' The dashboard only sends \codepath{true} after you press the explicit mutation action for operations such as rename, fill, drop, type conversion, or setting a cell.

The route clamps the row limit between 1 and 500. That keeps an accidental giant preview from overwhelming the page.

\section{How to add a CSV feature without hurting OTA}

Put new data operations in the Flask route or a Python helper. Add the matching control to \codepath{index.html} and its behavior to \codepath{app.js}. Keep its requests on \codepath{API_BASE}.

Likewise, an OTA feature belongs in Axum and \codepath{ota.js}. The two workbenches may share the front desk, but neither needs to borrow the other's tools.

\begin{warningbox}
Do not rename the existing CSV routes while adding firmware features. Saved frontend behavior already depends on those addresses. Small additions around the working flow are safer than rewriting it.
\end{warningbox}

\begin{checkpoint}
Profile the four-row file, remove its duplicate, fill the missing temperature, and download the result. Refresh the browser and recover the session from Previous Runs. Then explain why the report survives even though JavaScript memory was refreshed.
\end{checkpoint}



\part{The Rust OTA Service}
\partintro{Now we cross into the Rust server. We will unpack one firmware ZIP carefully and follow it until a real application image and its build history appear.}
\chapter{Meet the Rust OTA Server}

\section{Read main.rs as a morning routine}

Open \codepath{ota-server/src/main.rs}. The file may look formal at first, yet its job is close to opening a workshop in the morning. Turn on the lights. Unlock the cabinets. Put the tools on the bench. Open the front door.

The Rust version of that routine goes in this order:

\begin{enumerate}
  \item Start logging so the terminal can tell us what happens.
  \item Read settings from \codepath{.env}.
  \item Open the SQLite database.
  \item Prepare the firmware builder.
  \item Prepare the coordinator that watches each build.
  \item Put those shared pieces into \codepath{AppState}.
  \item Connect the HTTP routes.
  \item Listen on the configured address and port.
\end{enumerate}

Here is the middle of that routine with some details folded away:

\begin{lstlisting}[language=Rust]
let config = Arc::new(Config::load()?);
let storage = Storage::connect(&config.database_url()).await?;
let builder = Arc::new(LocalFirmwareBuilder::new(/* settings */));
let coordinator = Arc::new(BuildCoordinator::new(
    storage.clone(),
    builder,
    config.project_dir.clone(),
    config.build_dir.clone(),
));

let state = AppState { config, storage, coordinator };
\end{lstlisting}

The word \codepath{Arc} means several requests can safely hold a shared reference to the same service object. You do not need to master it now. Notice the shape: configuration, storage, and building meet in one state bundle, and route handlers receive that bundle when they need it.

\section{Configuration is the page of house rules}

Move to \codepath{ota-server/src/config.rs}. This file reads each \codepath{OTA_...} setting, supplies local defaults, checks values, and resolves folder paths from the repository root.

It also creates \codepath{firmware-projects}, \codepath{firmware-builds}, and \codepath{data} when they are missing. That is why a fresh copy of the project can start without you building those folders by hand.

If no public URL is supplied, the server tries to detect a likely LAN address. Detection is convenient on a simple home network. An explicit \codepath{OTA_PUBLIC_BASE_URL} is easier to reason about when a PC has Ethernet, Wi-Fi, virtual adapters, or a VPN.

When startup succeeds, the terminal prints two views:

\begin{lstlisting}[numbers=none]
OTA Server running
Local:
http://localhost:7000

LAN / public firmware URL:
http://192.168.1.5:7000
\end{lstlisting}

The local line is for the browser on the PC. The LAN line is the address an ESP32 can reach.

\section{Routes are the labels on the front desk}

Open \codepath{ota-server/src/routes/mod.rs}. This file is short because it mainly matches an address with the function that handles it.

The work is grouped into three nearby files:

\begin{description}
  \item[\codepath{routes/firmware.rs}] Receives project ZIP files and serves finished firmware.
  \item[\codepath{routes/builds.rs}] Starts builds and returns build history or logs.
  \item[\codepath{routes/devices.rs}] Remembers boards, heartbeats, deployments, update checks, and OTA progress.
\end{description}

Find this line:

\begin{lstlisting}[language=Rust,numbers=none]
.route("/api/ota/status", get(status))
\end{lstlisting}

Read it from left to right: when a GET request arrives at that address, call \codepath{status}. Other lines follow the same grammar.

\section{Why the browser is allowed through}

The page comes from port 5000 and calls a server on port 7000. Browsers treat those as different origins. Axum therefore needs a CORS rule that names the dashboard origins it trusts.

The default list contains:

\begin{lstlisting}[numbers=none]
http://localhost:5000
http://127.0.0.1:5000
\end{lstlisting}

If you serve the frontend from another address, add that exact origin to \codepath{OTA_ALLOWED_ORIGINS}. A CORS error in the browser can happen even while the Axum status URL opens normally in its own tab.

\begin{warningbox}
Keep the origin list narrow. Allowing every website to call a local firmware builder turns a browser convenience into a needless open door.
\end{warningbox}

\section{Errors should speak in complete sentences}

Open \codepath{ota-server/src/errors.rs}. Every ordinary API error carries an HTTP status, a short machine-readable code, and a message for a person. It may also carry details and a hint.

\begin{lstlisting}[language=JSON,numbers=none]
{
  "success": false,
  "error": {
    "code": "INVALID_FIRMWARE_PROJECT",
    "message": "Managed application entrypoint was not found",
    "details": "src/main.rs"
  }
}
\end{lstlisting}

The frontend can show the message without guessing what went wrong. A developer can use the short code in tests. The status number gives HTTP clients the broad category:

\begin{description}
  \item[400] The request could not be read.
  \item[404] The named record does not exist.
  \item[409] The request clashes with current state, such as a second active build.
  \item[413] The upload is too large.
  \item[422] The request was readable but failed a project rule.
  \item[500] Something inside the server failed unexpectedly.
\end{description}

\section{SQLite is the server's notebook}

The database file is \codepath{data/ota.db}. You do not need to create its tables. \codepath{storage.rs} does that during startup.

Each table answers a different question:

\begin{longtable}{p{0.19\textwidth}p{0.7\textwidth}}
\toprule
Table & The question it answers\\
\midrule
\endhead
projects & Which ZIP was uploaded, for which target and version?\\
builds & What happened when we tried to compile it?\\
firmware & Which successful binary can a device download?\\
devices & Which boards have introduced themselves, and what did they report?\\
deployments & Which firmware was assigned to which board, and how far did it get?\\
\bottomrule
\end{longtable}

SQLite foreign keys tie those pages together. A firmware file that still belongs to a device or deployment cannot be casually deleted. The API returns a conflict instead of tearing a page out of the notebook and leaving dangling references behind.

\begin{checkpoint}
Open \codepath{main.rs}, \codepath{config.rs}, \codepath{routes/mod.rs}, and \codepath{storage.rs} in four tabs. In each tab, find the one place that matches its chapter role: startup, settings, addresses, and database tables.
\end{checkpoint}


\chapter{What Really Happens to the ZIP}

\section{Begin at the upload form}

On the OTA Firmware page, the form asks for three things: a ZIP, a board, and a version. The browser packs them into a multipart request with these field names:

\begin{center}
\begin{tabularx}{0.9\textwidth}{l l Y}
\toprule
Field & Example & Why the server needs it\\
\midrule
\codepath{project} & \codepath{sensor.zip} & The Rust source and Cargo files\\
\codepath{target} & \codepath{esp32s3} & The compiler and chip choice\\
\codepath{version} & \codepath{1.0.1} & The release name reported by the new firmware\\
\bottomrule
\end{tabularx}
\end{center}

Axum gives the upload a random project ID. That ID names its private extraction folder. A filename supplied by the browser never gets to choose where the server writes.

\section{The server checks the suitcase before unpacking it}

A ZIP is much like a suitcase with a list of intended destinations. A dishonest entry might claim that its file belongs in \codepath{../../somewhere-else}. Extracting that name without checking would let the archive escape its room. This attack is often called ZIP slip.

The extractor refuses parent paths, absolute paths, symbolic links, more than 10,000 entries, and more extracted data than the configured limit. It also insists on finding a \codepath{Cargo.toml}.

These checks may feel strict when your own archive is rejected. They protect every other folder on the PC. Fix the ZIP layout rather than relaxing the guard.

\section{The project chooses one of two paths}

After extraction, \codepath{services/project.rs} looks for \codepath{ota-app.toml}.

If the file exists, the project is a \emph{managed application}. It supplies the product behavior while the stable runtime supplies Wi-Fi, OTA, and rollback.

\begin{lstlisting}[language=TOML,numbers=none]
schema_version = 1
application_api_version = 2
entrypoint = "src/main.rs"
\end{lstlisting}

If the file is absent, the project is \emph{standalone firmware}. Axum builds the project as it arrived. That project must carry its own OTA client if it should update itself again.

For a first application, managed mode is easier. You can write sensor behavior without copying the network and partition code into every upload.

\section{How two files named main.rs live together}

This point deserves a slow pass.

The stable runtime has a normal executable \codepath{fn main()} in \codepath{examples/esp32-rust-ota-client/src/main.rs}. It owns boot, Wi-Fi, OTA, and application startup.

Your ZIP also keeps \codepath{src/main.rs}, shaped like an ordinary firmware program. It exports two free functions:

\begin{lstlisting}[language=Rust]
pub fn setup(context: &mut AppContext) -> anyhow::Result<State>;
pub fn main(context: AppContext, state: State) -> anyhow::Result<()>;
\end{lstlisting}

Axum copies the stable runtime into a private generated workspace. It places your application under \codepath{applications/device-app} and adjusts only the copied Cargo manifests. In that copy, Cargo treats your \codepath{src/main.rs} as a library source file. The runtime can then call \codepath{device_app::setup} and \codepath{device_app::main}.

Your uploaded files stay unchanged. Rust links both pieces into one complete image.

\begin{mentorbox}[A kitchen picture]
The stable runtime is a kitchen with power, water, and safety rules. Your upload is the recipe being cooked today. The finished meal includes both. Uploading a new recipe does not pile a second kitchen on top of the first one.
\end{mentorbox}

\section{The builder checks its tools}

The actual command runner lives in \codepath{services/builder.rs}. Before it touches the project, it asks three programs for their versions:

\begin{lstlisting}
cargo --version
rustc --version
espflash --version
\end{lstlisting}

A missing tool becomes a failed build record with a clear message. Axum stays alive, so you can fix the environment and try another build.

The builder is hidden behind a small Rust trait named \codepath{FirmwareBuilder}. Today the implementation runs tools directly on the PC. Tomorrow another implementation could run them inside a locked-down container while the routes remain the same.

\section{Cargo turns source into an ELF}

The build command is:

\begin{lstlisting}
cargo build --release --target xtensa-esp32s3-espidf
\end{lstlisting}

Axum does not build a shell sentence from uploaded text. It starts \codepath{cargo} and passes each argument separately. The working directory is the trusted project workspace.

While Cargo works, the server reads both output streams line by line and saves them in the build record. This is why the dashboard can show the real compiler message and why an old failed build still has useful logs after a browser refresh.

\section{espflash prepares the OTA image}

Cargo's ELF is useful to development tools, but the inactive OTA partition needs an ESP-IDF application image. The next command follows this shape:

\begin{lstlisting}
espflash save-image --chip esp32s3 --format esp-idf \
  <ELF_FILE> firmware-sensor-v1.0.1-a8c4f2.bin
\end{lstlisting}

The builder deliberately asks for an application image. It does not make a merged factory image containing bootloader and partition data. Those larger pieces belong to initial USB provisioning.

\section{The fingerprint closes the package}

The server reads the finished \codepath{.bin} and calculates SHA-256. Think of the hash as a long fingerprint. The device calculates the same fingerprint while it downloads. If the two strings differ, the bytes changed or the wrong file arrived, so the device refuses to activate it.

On success, SQLite receives the filename, byte size, target, version, project ID, build ID, timestamp, and hash. The dashboard then adds the image to the firmware table.

\begin{mentorbox}[A fingerprint is not a signature]
SHA-256 catches corruption. It does not tell the device who approved the firmware. Production publisher identity needs real signing and protected keys. The project leaves that as a visible future boundary.
\end{mentorbox}

\section{Read the build states as a little story}

\begin{center}
\ttfamily uploaded $\rightarrow$ queued $\rightarrow$ building $\rightarrow$ success\par
\hspace*{7.4cm}$\searrow$ failed
\end{center}

Uploaded means the suitcase passed inspection. Queued means a build record exists. Building means the tools are working. Success means the binary and metadata are real. Failed means the log holds the reason the journey stopped.

Only one queued or active build is allowed for the same project. A second click receives a conflict instead of launching two compilers into the same working folders.

\begin{warningbox}[Your own ZIP can still run code on the PC]
Cargo may execute \codepath{build.rs}, procedural macros, and dependency build scripts. Safe extraction protects paths, not the act of compilation. Use trusted projects until the builder runs inside a proper sandbox.
\end{warningbox}

\begin{checkpoint}
Open \codepath{examples/uploaded-firmware-basic}. Check that its ZIP root would contain \codepath{Cargo.toml}, \codepath{ota-app.toml}, and \codepath{src/main.rs}. Upload it as version \codepath{1.0.1} and watch the log. A successful toolchain should produce a real binary. A missing tool should produce a real failure. Both outcomes tell the truth.
\end{checkpoint}



\part{Firmware on the ESP32-S3}
\partintro{The ESP32 finally joins us. First it receives a dependable foundation over USB. After that, your smaller application code can ride inside each complete OTA release.}
\chapter{The Program That Lives on the ESP32}

\section{Give the board a dependable base}

Take the ESP32-S3 in your hand. Before OTA can help us, this board needs a program that knows how to join Wi-Fi, find the PC, and write a new image safely. That first program is the stable runtime under \codepath{examples/esp32-rust-ota-client}.

The word \emph{stable} does not mean it can never change. It means ordinary product updates should not have to rewrite its plumbing. It looks after:

\begin{itemize}
  \item ESP-IDF startup and logs;
  \item the Wi-Fi connection;
  \item device identity and heartbeats;
  \item update checks and downloads;
  \item the two OTA application slots;
  \item hash verification, reboot, and rollback;
  \item the moment your uploaded application begins running.
\end{itemize}

Your application can then spend its attention on a sensor, motor, display, or other product job.

\section{The first journey goes through USB}

Open PowerShell in the runtime directory. Fill in values for the network and the reachable server address:

\begin{lstlisting}
cd examples\esp32-rust-ota-client
$env:WIFI_SSID="your-ssid"
$env:WIFI_PASS="your-password"
$env:OTA_SERVER_URL="http://192.168.1.5:7000"
$env:DEVICE_NAME="Sensor node"
\end{lstlisting}

Check the IP carefully. It belongs to the PC running Axum.

When the board, cable, and partition table are ready, this command builds, flashes, and opens the monitor:

\begin{lstlisting}
cargo +esp run --release --target xtensa-esp32s3-espidf
\end{lstlisting}

The configured Cargo runner calls \texttt{espflash flash --monitor}. The first flash can install the bootloader and partition layout as well as the application.

\begin{warningbox}[Stop and identify the board]
Flashing writes to physical hardware. Confirm that the connected board is an ESP32-S3, the USB cable carries data, the flash size matches the partition plan, and the selected serial port belongs to this board.
\end{warningbox}

When the runtime boots, the serial terminal should show Wi-Fi activity. After it receives an IP address, it registers with Axum. Refresh Devices in the dashboard and look for the name you supplied.

\section{Read the true fn main like a story}

Open \codepath{examples/esp32-rust-ota-client/src/main.rs}. The function begins here:

\begin{lstlisting}[language=Rust,numbers=none]
fn main() -> Result<()> {
    esp_idf_svc::sys::link_patches();
    EspLogger::initialize_default();
    // platform setup continues...
}
\end{lstlisting}

Continue down the file slowly. The runtime takes the board's peripherals, separates the modem, starts the system event loop and NVS, then constructs \codepath{WifiManager}. It derives a device ID from the station MAC unless you supplied one.

Next it builds \codepath{AppContext}. This is the bag handed to your application. Before calling your application, it asks ESP-IDF whether the current image is new and still waiting for validation.

Then comes the handoff:

\begin{lstlisting}[language=Rust]
let application_state = device_app::setup(&mut context)?;

if unverified_image {
    mark_running_image_valid()?;
}

// Start the OTA supervisor in its own thread.
device_app::main(context, application_state)
\end{lstlisting}

The real source contains fuller error handling. The shape above is the part to remember: prepare, test application setup, approve the new slot, keep OTA alive, and enter the application's main loop.

\section{Open the smaller crates}

The runtime is divided so each piece can be held in your head:

\begin{description}
  \item[\codepath{crates/application-api}] Defines what the application is allowed to receive.
  \item[\codepath{crates/wifi-manager}] Owns the modem and keeps trying when the network drops.
  \item[\codepath{crates/ota-runtime}] Talks to Axum and writes an update into the inactive slot.
  \item[\codepath{applications/default-app}] Supplies the small behavior used during the first USB flash.
\end{description}

You can open one crate without reading the others. This is especially useful when a Wi-Fi problem and a sensor problem happen on the same day.

\section{AppContext is the handover bag}

The application receives this broad shape:

\begin{lstlisting}[language=Rust]
pub struct AppContext {
    pub device: DeviceInfo,
    pub network: NetworkHandle,
    pub peripherals: ApplicationPeripherals,
}
\end{lstlisting}

\codepath{device} tells the application its ID, friendly name, chip, and firmware version. \codepath{network} lets it ask whether Wi-Fi is connected and which IP was last reported. \codepath{peripherals} carries the available pins and controllers.

Each hardware item is wrapped in \codepath{Option}. When your setup function takes I2C0, that option becomes empty. The code now has a visible record that one part of the application owns that controller.

The modem is missing from the bag on purpose. \codepath{wifi-manager} keeps it so heartbeat and OTA polling can continue while your application runs.

\begin{mentorbox}[Would a Wi-Fi manager become a problem?]
It becomes helpful when ownership stays in one place. The current manager already owns the modem, event loop, NVS-backed Wi-Fi service, and reconnect loop. A future setup access point or captive portal can grow there. Uploaded applications can keep reading the same network handle.
\end{mentorbox}

\section{A new flash does not stack old code}

Imagine two labeled drawers, \codepath{ota_0} and \codepath{ota_1}. The board runs one drawer while the updater writes a complete image into the other. After verification and reboot, the drawers trade roles.

The newly written image contains the runtime and the chosen application together. It replaces the inactive slot from beginning to end. Old structs, functions, and libraries do not remain loaded beside the new program.

The previous active drawer stays available for rollback until the new image proves it can start. On the next release, ESP-IDF may overwrite that older drawer. The PC can keep a long firmware history in SQLite, while the board still has only its two application slots.

\begin{checkpoint}
Before connecting hardware, find the runtime's \texttt{fn main}, the call to \codepath{device_app::setup}, the call to \codepath{device_app::main}, and the place where the modem is separated from application peripherals. Trace them in that order with your finger.
\end{checkpoint}

\chapter{Write Your First Uploaded Firmware}

\section{Copy the small example before starting from nothing}

Open \codepath{examples/uploaded-firmware-basic}. It contains only the pieces the server expects:

\begin{lstlisting}[numbers=none]
Cargo.toml
ota-app.toml
src/
  main.rs
\end{lstlisting}

Copy this folder somewhere safe and give the Cargo package a new name. Beginning from a working example lets you change one idea at a time.

The descriptor tells Axum how to treat the project:

\begin{lstlisting}[language=TOML,numbers=none]
schema_version = 1
application_api_version = 2
entrypoint = "src/main.rs"
\end{lstlisting}

Leave these values as they are for your first application.

\section{Your program still feels like main.rs}

Open \codepath{src/main.rs}. There is no application class and no framework trait to implement. You write two public functions.

\begin{lstlisting}[language=Rust]
use anyhow::Result;
use firmware_app_api::AppContext;

pub fn setup(context: &mut AppContext) -> Result<u64> {
    log::info!("Hello from {}", context.device.device_id);
    Ok(0)
}

pub fn main(context: AppContext, mut count: u64) -> Result<()> {
    loop {
        count += 1;
        log::info!(
            "count={count}, wifi={}",
            context.network.is_connected()
        );
        std::thread::sleep(std::time::Duration::from_secs(10));
    }
}
\end{lstlisting}

\codepath{setup} runs once after the runtime has prepared the board. \codepath{main} is your long-running loop. The stable entrypoint calls both for you.

Change the hello message. That one line is enough for a first personalized build.

\section{The value between setup and main is your toolbox}

The example returns a number from setup and receives it as \codepath{count} in main. That number is the application state.

Your own state can be as small as \codepath{()} when nothing needs to be carried. It can also be a struct containing a sensor driver, counters, buffers, or display state.

\begin{lstlisting}[language=Rust]
pub struct State {
    sample_number: u64,
    alarm_enabled: bool,
}

pub fn setup(context: &mut AppContext) -> Result<State> {
    log::info!("Preparing peripherals for {}", context.device.name);
    Ok(State {
        sample_number: 0,
        alarm_enabled: true,
    })
}

pub fn main(context: AppContext, mut state: State) -> Result<()> {
    loop {
        state.sample_number += 1;
        // Read the sensor and do your real work here.
    }
}
\end{lstlisting}

This is ordinary Rust data, not a class required by the OTA system.

\section{Take a peripheral during setup}

Suppose the application needs I2C0 and some pins. The beginning may look like this:

\begin{lstlisting}[language=Rust]
use anyhow::{Context, Result};

pub fn setup(context: &mut AppContext) -> Result<State> {
    let pins = context.peripherals.pins.take()
        .context("GPIO pins were already taken")?;
    let i2c0 = context.peripherals.i2c0.take()
        .context("I2C0 was already taken")?;

    // Construct the ESP-IDF I2C driver and sensor here.
    // Return the finished driver inside State.
}
\end{lstlisting}

The calls to \codepath{take()} move those resources into your code. If setup is called with a missing resource, it returns a readable error. On a brand-new OTA image, that error can trigger rollback.

Keep fallible driver construction and a short startup check in setup. Do not let setup loop forever. The runtime needs a clear success or failure before it can approve the new slot.

\begin{warningbox}
Return owned state from setup. A reference borrowed from \codepath{AppContext} cannot safely outlive the mutable setup call when the context later moves into main. Real ESP-IDF drivers may need careful lifetime and type aliases; follow the driver crate's example for your exact chip and pins.
\end{warningbox}

\section{Add ordinary Cargo libraries}

Yes, your uploaded application can use libraries. Open \codepath{examples/uploaded-firmware-with-libs/Cargo.toml}:

\begin{lstlisting}[language=TOML,numbers=none]
[dependencies]
anyhow = "1"
firmware-app-api = "0.1"
heapless = "0.8"
log = "0.4"
serde = { version = "1", features = ["derive"] }
serde_json = "1"
\end{lstlisting}

That example keeps samples in a fixed-capacity queue and turns telemetry into JSON. Sensor drivers, \codepath{embedded-hal} helpers, compact collections, and serialization crates can work the same way.

The library must support \codepath{xtensa-esp32s3-espidf}. A desktop crate that assumes a mouse, filesystem, or full operating system may compile poorly or make little sense on the board. Check the crate's examples and features.

If you use a local path dependency, include its whole folder inside the ZIP. The path must stay under the uploaded project root. Axum rejects paths that reach outside the archive.

\begin{mentorbox}[Why firmware-app-api looks like a normal dependency]
The uploaded manifest uses \codepath{firmware-app-api = "0.1"} as a portable placeholder. In its private build copy, Axum points that name to the matching API crate inside the stable runtime. Your original manifest is left alone.
\end{mentorbox}

\section{Give the release a useful version}

The version typed into the dashboard becomes the firmware version compiled into the managed runtime and stored in SQLite. Begin with a calm sequence such as \codepath{1.0.0}, \codepath{1.0.1}, and \codepath{1.1.0}.

Deployment selects a specific firmware ID. The current server checks whether the desired version equals the device's current version; it does not sort semantic versions and choose one by itself.

\section{Zip the contents, not the outer shelf}

On Windows PowerShell:

\begin{lstlisting}
Compress-Archive \
  -Path examples\uploaded-firmware-basic\* \
  -DestinationPath uploaded-firmware-basic.zip
\end{lstlisting}

On Linux or macOS:

\begin{lstlisting}
cd examples/uploaded-firmware-basic
zip -r ../../uploaded-firmware-basic.zip \
  Cargo.toml ota-app.toml src README.md
\end{lstlisting}

Open the ZIP and check its first level. You should immediately see \codepath{Cargo.toml}, \codepath{ota-app.toml}, and \codepath{src}. Upload it on OTA Firmware. The device will receive the later \codepath{.bin}, never this ZIP.

\begin{checkpoint}
Make a copy of the basic example. Change its package name and hello message. Let setup return a small state struct of your own. Zip the contents and upload version \codepath{1.0.1}. The upload response should say \codepath{managed_application}, API version 2, and \codepath{src/main.rs}.
\end{checkpoint}


\chapter{Watch an OTA Update Travel}

\section{The dashboard leaves a note for the device}

Pressing Deploy OTA does not force a connection into the ESP32. The dashboard tells Axum which firmware the device should want. Axum records that choice as pending.

The ESP32 checks the server on its own schedule. On the next check, Axum hands back the note. This is a pull-based update: the device opens every network conversation.

\begin{center}
\begin{tcolorbox}[width=0.92\textwidth,colback=BookPaper,colframe=BookTealLight]
\centering\ttfamily
You press Deploy OTA\par
$\downarrow$\par
Axum records the desired firmware\par
$\downarrow$ the ESP32 asks for an update\par
download $\rightarrow$ check $\rightarrow$ install $\rightarrow$ reboot\par
$\downarrow$\par
the new runtime reports its version
\end{tcolorbox}
\end{center}

A sleeping or temporarily offline board can catch the note later. The PC does not need to know how to open a connection through the board's firewall or changing address.

\section{The device introduces itself}

Once Wi-Fi is connected, the runtime posts a registration like this:

\begin{lstlisting}[language=JSON,numbers=none]
{
  "device_id": "ESP32-S3-A1B2C3D4E5F6",
  "name": "Sensor node",
  "chip": "esp32s3",
  "current_version": "1.0.0",
  "ip_address": "192.168.1.24"
}
\end{lstlisting}

That request creates the card on the Devices page. The server accepts a tidy device ID made from letters, numbers, dots, dashes, or underscores.

The runtime then sends a heartbeat during each cycle. A heartbeat is a short ``I am still here'' message with the current version, status, and IP. Axum updates \codepath{last_seen} each time.

An old card can remain after a board loses power. Look at \codepath{last_seen} before trusting the word Online.

\section{The update check returns a small manifest}

When nothing is assigned, the device receives:

\begin{lstlisting}[language=JSON,numbers=none]
{"update_available": false}
\end{lstlisting}

When version \codepath{1.0.1} is waiting, it receives a description of the exact binary:

\begin{lstlisting}[language=JSON,numbers=none]
{
  "update_available": true,
  "firmware_id": "c53b...",
  "version": "1.0.1",
  "target": "esp32s3",
  "size": 945332,
  "sha256": "f2534c...",
  "download_url": "http://192.168.1.5:7000/api/ota/firmware/c53b.../download"
}
\end{lstlisting}

The device checks the target, size, hash shape, and URL before opening an OTA writer. The server also prevents you from assigning ESP32-S3 firmware to a differently registered chip.

\section{The binary moves one cup at a time}

A firmware image can be much larger than the spare RAM on a small device. The client therefore uses a 4 KiB buffer.

Picture moving water from a barrel with one cup. Read a cup from HTTP. Pour it into the inactive flash slot. Add those same bytes to the running SHA-256 calculation. Repeat until the response ends. The whole barrel never has to sit in memory.

As bytes arrive, the runtime reports downloading progress in roughly five-percent steps. At the end it checks the total byte count and compares the finished fingerprint with the server's value.

If size or hash differs, the writer is aborted. The current running application remains untouched.

\section{The device narrates the update}

Status callbacks let the dashboard follow along:

\begin{lstlisting}[language=JSON,numbers=none]
{"status":"downloading","progress":35}
{"status":"verifying","progress":100}
{"status":"installing","progress":100}
{"status":"rebooting","progress":100}
\end{lstlisting}

A failure carries the reason:

\begin{lstlisting}[language=JSON,numbers=none]
{
  "status": "failed",
  "error": "SHA256 mismatch"
}
\end{lstlisting}

After reboot, the new image runs setup. If that succeeds, the runtime marks the slot valid and reports:

\begin{lstlisting}[language=JSON,numbers=none]
{
  "status": "success",
  "progress": 100,
  "current_version": "1.0.1"
}
\end{lstlisting}

Axum clears the desired firmware. The device will not be offered the same release on every poll.

\section{Look inside the two-drawer partition table}

The supplied 4 MiB example is:

\begin{lstlisting}[numbers=none]
# Name,   Type, SubType, Offset,   Size
nvs,      data, nvs,     0x9000,   0x6000
otadata,  data, ota,     0xf000,   0x2000
phy_init, data, phy,     0x11000,  0x1000
ota_0,    app,  ota_0,   0x20000,  0x1e0000
ota_1,    app,  ota_1,   0x200000, 0x1e0000
\end{lstlisting}

\codepath{otadata} remembers which application slot should boot. \codepath{ota_0} and \codepath{ota_1} are the two drawers. ESP-IDF chooses the inactive one; application code should not assume a fixed winner.

\section{Rollback gives a new image a trial period}

The example enables \codepath{CONFIG_BOOTLOADER_APP_ROLLBACK_ENABLE=y}. When an OTA image boots for the first time, its slot is unverified.

The runtime calls your setup function. A successful setup marks the slot valid. A failed setup marks it invalid and asks for a reboot. The rollback-aware bootloader can then return to the previous valid slot.

This safety net depends on the partition table, bootloader configuration, and application validation call working together. Test it with a deliberately failing setup on your own hardware before relying on it.

\begin{warningbox}
SHA-256 and plain HTTP are suitable for corruption checks on a trusted development LAN. A production device also needs an authenticity plan such as signed images, protected keys, Secure Boot, and an intentional transport policy.
\end{warningbox}

\begin{checkpoint}
Tell the update story without endpoint names: leave a note, let the board ask, move the image in small cups, compare fingerprints, switch drawers, test startup, and report the new version. Then match each step to the code or dashboard screen.
\end{checkpoint}



\part{Practice and Extension}
\partintro{The last part is a workbench. You will run the whole path, learn how to read common failures, and sketch improvements without pulling apart the pieces that already work.}
\chapter{Put Everything Together}

\section{Clear a little space on the desk}

This is the day we run the whole path. Keep the PC, ESP32-S3, data USB cable, and Wi-Fi details nearby. Open one terminal for Flask and one for Axum. Leave room for a third terminal when the serial monitor appears.

There is no prize for rushing. At the end of each stage, look for the small piece of evidence named on the page. If it is missing, stay with that stage.

\section{Stage one: wake both PC services}

Open \codepath{.env} and check the PC's LAN address, Wi-Fi name, Wi-Fi password, and device name. Then start the project:

\begin{lstlisting}[numbers=none]
# Windows
.\run-dev.bat

# Linux or macOS
./run-dev.sh
\end{lstlisting}

Open these two addresses:

\begin{lstlisting}[numbers=none]
http://localhost:5000
http://localhost:7000/api/ota/status
\end{lstlisting}

\textbf{What to notice:} the first tab shows the dashboard; the second says the OTA service is online and prints the correct LAN-facing public URL.

\section{Stage two: give the old CSV path a quick health check}

Upload the small \codepath{sensor-readings.csv} from Chapter 4. Watch progress, open the report, and run a playground preview. Visit Previous Runs and make sure the session is listed.

\textbf{What to notice:} switching to OTA Firmware and back does not lose the CSV session. This proves our newer pages still respect the original Flask work.

\section{Stage three: pack the basic application}

On Windows PowerShell:

\begin{lstlisting}
Compress-Archive \
  -Path examples\uploaded-firmware-basic\* \
  -DestinationPath uploaded-firmware-basic.zip
\end{lstlisting}

Open the ZIP. Its first level should show \codepath{Cargo.toml}, \codepath{ota-app.toml}, and \codepath{src}.

Go to OTA Firmware. Choose the ZIP, select ESP32-S3, type version \codepath{1.0.1}, and press Upload \& Build.

\textbf{What to notice:} the badge moves from uploading to queued and building. The terminal begins with tool versions and the managed-application mode. Real compiler lines follow.

If the build fails, keep the record. Find the first useful error and fix that condition. Success should never be invented by editing SQLite or dropping a random binary into \codepath{firmware-builds}.

\section{Stage four: inspect what a successful build made}

When the badge reaches success, look at the firmware table. The new row should have version \codepath{1.0.1}, target ESP32-S3, a filename, byte size, SHA-256, time, and Ready status.

Behind the screen, you will find:

\begin{itemize}
  \item the extracted upload under \codepath{firmware-projects};
  \item a generated managed workspace inside that project;
  \item the application binary under \codepath{firmware-builds};
  \item matching project, build, and firmware rows in \codepath{data/ota.db}.
\end{itemize}

Press Download once. The downloaded byte count should match the row in the table.

\section{Stage five: give the board its first base program}

This is the physical boundary. Review \codepath{examples/esp32-rust-ota-client/partitions.csv} and compare it with the board's flash size. Connect the board with a data cable.

Set the runtime values and flash:

\begin{lstlisting}
cd examples\esp32-rust-ota-client
$env:WIFI_SSID="your-ssid"
$env:WIFI_PASS="your-password"
$env:OTA_SERVER_URL="http://192.168.1.5:7000"
$env:DEVICE_NAME="Tutorial sensor"
cargo +esp run --release --target xtensa-esp32s3-espidf
\end{lstlisting}

Replace the sample IP before pressing Enter.

\begin{warningbox}
The command writes to the connected board. Confirm the chip, serial port, flash size, partition table, and stable power first. Do not reuse an example partition map on an existing product until its data partitions have been reconciled.
\end{warningbox}

\textbf{What to notice:} the serial terminal shows startup and Wi-Fi connection. Soon afterward, the board appears on the Devices page with its current version and IP address.

\section{Stage six: leave the update note}

Open Devices and press Deploy OTA for the new board. Choose the firmware row for \codepath{1.0.1}.

The card should first say Pending. The runtime checks every thirty seconds, so give it a moment. Then watch Downloading, Verifying, Installing, and Rebooting.

The serial monitor may disappear briefly during restart. Open it again if needed. The new image runs application setup, marks its slot valid, and starts the uploaded main loop.

\textbf{What to notice:} the device card reaches Success and its current version becomes \codepath{1.0.1}. The serial output contains the basic application's cycle message.

\section{Stage seven: change something you can recognize}

Copy the basic application. Change its log line to a message you will recognize immediately, such as the room or sensor name. Build it as \codepath{1.0.2}, then deploy it.

After reboot, look for the new message in the serial monitor. This is the satisfying moment: source from your ZIP is now running inside the complete firmware image built by the PC.

If you want to try dependencies next, use \codepath{examples/uploaded-firmware-with-libs}. It already demonstrates \codepath{heapless}, Serde, and JSON.

\section{Write down the evidence while it is fresh}

A short lab note can save an afternoon later. Record the build ID, firmware ID, version, byte size, hash, device ID, old and new versions, board model, tool versions, and final serial message. Add the real error and solution when a stage needed repair.

\begin{checkpoint}[The whole path]
You have completed the lab when the CSV report still works, a real toolchain build creates firmware, a physically provisioned device registers, deployment becomes an update manifest, the device streams and verifies the binary, and the rebooted image reports version \codepath{1.0.1} or later.
\end{checkpoint}


\chapter{When Something Goes Wrong}

\section{Sit with the last thing that worked}

An error becomes easier when you give it an address. Ask one question: what was the last step that definitely worked?

If the dashboard never opened, stay with Flask. If CSV works and the OTA light is dark, stay with Axum. If Axum accepted the ZIP and Cargo failed, stay with the build log. If the PC can download the binary and the board cannot, stay with the network between them.

This habit is slower for ten seconds and faster for the next hour.

\section{The dashboard page will not open}

Look at the Flask terminal. If it names a missing package, install \codepath{python/requirements.txt} into \codepath{python/.venv}. Then try:

\begin{lstlisting}[numbers=none]
http://127.0.0.1:5000/
\end{lstlisting}

If that address opens, Flask is serving the page. If the page has no style or buttons do nothing, open \codepath{/style.css} and \codepath{/app.js} on the same port to see whether those files are served.

Do not open \codepath{index.html} by double-clicking it. The resulting \codepath{file:///} page lives in a different browser world from the Flask-served dashboard.

\section{CSV works and OTA says offline}

Open:

\begin{lstlisting}[numbers=none]
http://localhost:7000/api/ota/status
\end{lstlisting}

If it fails, read the Axum terminal. Cargo may still be compiling. Configuration may contain a bad port or URL. Another program may already own port 7000.

If the status opens in its own tab but the dashboard request is blocked, open the browser's Network panel and look for CORS. Compare the dashboard's exact origin with \codepath{OTA_ALLOWED_ORIGINS}.

\section{The ZIP is turned away}

Read the error body before repacking anything. It usually points to one of these:

\begin{itemize}
  \item the chosen file is not a ZIP;
  \item \codepath{Cargo.toml} is missing;
  \item \codepath{ota-app.toml} uses an old API version;
  \item \codepath{src/main.rs} is missing or the entrypoint points outside the project;
  \item a local dependency reaches outside the archive;
  \item the archive contains a symlink or unsafe path;
  \item the encoded or extracted size is too large.
\end{itemize}

Open the ZIP and look at its first level. The clean managed layout is \codepath{Cargo.toml}, \codepath{ota-app.toml}, and \codepath{src}.

\section{The build stops before your code appears}

The first log lines test \codepath{cargo}, \codepath{rustc}, and \codepath{espflash}. If one is missing, run its \codepath{--version} command in the same terminal environment used by Axum.

After installing a tool, close and restart the OTA terminal. A running process keeps the environment it had when it started.

On Windows, \codepath{link.exe not found} points to Visual Studio Build Tools with Desktop development with C++. The Espressif toolchain alone does not supply that host linker.

\section{Rust gives you a page of compiler errors}

Scroll to the first clear error with a source path and line number. Later errors often grow from that first one.

For an uploaded application, check the basic handshake:

\begin{lstlisting}[language=Rust,numbers=none]
pub fn setup(context: &mut AppContext) -> anyhow::Result<State>
pub fn main(context: AppContext, state: State) -> anyhow::Result<()>
\end{lstlisting}

The state type returned by setup must match the second main parameter. Take hardware from \codepath{context.peripherals}. Check that new libraries support the ESP32-S3 ESP-IDF target.

Make one correction and build again. Changing five things at once can hide which one mattered.

\section{Cargo succeeds and image creation fails}

Find the logged ELF path. The next lines should show the \codepath{espflash save-image} stage. Run \codepath{espflash --version} in the Axum terminal and confirm the expected command is available.

The server will only accept a non-empty application image in \codepath{firmware-builds}. A successful Cargo line does not excuse a missing \codepath{.bin}.

\section{The PC sees firmware and the board does not}

Look at \codepath{public_base_url} in the status response. It must contain the PC's LAN IP. Then try that LAN URL from a phone or another computer on the same Wi-Fi.

If the second device cannot open it, check the PC firewall, network profile, VPN routes, and access-point client isolation. If the second device can open it, compare the exact URL compiled into the ESP32 runtime.

\section{The device card never appears}

The dashboard does not discover USB devices. Registration must come from the running ESP32 program.

Open the serial monitor. Look first for Wi-Fi connection. Then look for the registration POST. Confirm that the board received the right SSID, password, and LAN server URL during its build.

If Wi-Fi provisioning has never happened, return to the first USB flash in Chapter 7.

\section{The card stays on Pending}

Pending means Axum saved the deployment. It has not heard completion from the device.

Check \codepath{last_seen}. The runtime normally polls every thirty seconds. Make sure the desired version differs from the current version and the device chip matches the firmware target. The serial terminal should show the next update check or its network error.

\section{Size or SHA-256 does not match}

The runtime aborts the inactive writer. That is the safe response. A size mismatch suggests a truncated download or wrong metadata. A hash mismatch means the received bytes do not match the artifact Axum recorded.

Download the same firmware on the PC, inspect the file size, and compare its hash with the metadata. Do not remove the device-side check to make the update continue.

\section{The binary has outgrown its drawer}

The sample OTA slots are each \codepath{0x1e0000} bytes. Compare the binary size with the real partition size. Reduce code or dependencies first. A larger slot requires a carefully revised partition table and usually a new serial provisioning step.

\section{Delete answers with 409 Conflict}

The firmware still belongs to a device's desired state or deployment history. The refusal protects those links. Do not delete the binary by hand while SQLite still points to it.

\section{Where are the ESP32 terminal messages?}

They are still in the serial monitor. Build History shows PC compiler output. Devices shows heartbeats and OTA status. Current dashboard code does not stream board logs.

Use \codepath{espflash monitor} today. The next chapter shows how a future serial bridge or network logger could bring those messages into a clearly labeled dashboard panel.

\begin{checkpoint}
Take one real failure from your terminal. On paper, write the last successful step, the failing boundary, the smallest test for that boundary, and one change to try. Keep the original error beside the result.
\end{checkpoint}


\chapter{Where You Can Take It Next}

\section{Choose one seam and keep the rest quiet}

The project is already divided along useful seams. A good extension changes the smallest responsible piece and lets its neighbors keep their familiar contracts.

Before writing code, finish this sentence: ``The new behavior belongs to \underline{\hspace{4cm}} because that piece already owns \underline{\hspace{4cm}}.''

If the answer is Wi-Fi credentials, open \codepath{wifi-manager}. If it is a new HTTP address, open \codepath{routes}. If it is a new build environment, open the \codepath{FirmwareBuilder} implementation.

\section{Welcome another ESP32 chip}

The current board list contains ESP32-S3. A second target needs more than another dropdown option.

Add its accepted name in target validation. Give \codepath{TargetSpec} the proper Rust target and \codepath{espflash} chip. Add the board choice in the frontend. Then make a runtime peripheral split that matches the actual chip, partition example, and rollback configuration.

An ESP32-C3, for example, uses RISC-V and does not offer the same hardware bundle as the S3. The application API should promise only resources the board truly has.

Begin with a known-good empty application. Prove build and serial flash before testing OTA and rollback.

\section{Turn Wi-Fi setup into a friendly first-run flow}

The current manager receives credentials at build time and reconnects when needed. A future board could behave like a new home device:

\begin{enumerate}
  \item Look for saved credentials in NVS.
  \item Try the saved network for a limited time.
  \item Start its own temporary setup access point when no network works.
  \item Show a small page where the owner chooses Wi-Fi.
  \item Save the credentials and return to station mode.
  \item Publish each change through the existing network handle.
\end{enumerate}

Keep the modem, event loop, NVS, and Wi-Fi service inside \codepath{wifi-manager}. The sensor application should not have to learn captive-portal code.

\section{Bring the ESP32 terminal into the dashboard}

There are two useful roads, and they show different parts of the journey.

\subsection{Road one: a serial bridge on the PC}

A small local service opens a selected COM or tty port, reads lines, and streams them to the browser with SSE or WebSocket. It can show bootloader messages, panics, and Wi-Fi failures before the network comes alive.

The dashboard must ask the user which port to open. Only one reader should own that port. The bridge needs bounded memory, clean disconnects, and a visible warning when flashing cannot begin because monitoring still owns the serial connection.

\subsection{Road two: logs sent over Wi-Fi}

The runtime can gather structured log records, post small batches to Axum, and let the browser subscribe to a device stream. This works after the USB cable is gone. It cannot report a failure that prevents Wi-Fi from starting.

A good dashboard labels each line as Serial or Network. The server also needs retention limits so a noisy loop cannot fill SQLite or memory.

\begin{mentorbox}[What exists today]
Neither device-log road is implemented in the current API. The existing dark terminal is only for PC build logs. Keeping that sentence visible prevents a mock terminal from being mistaken for hardware evidence.
\end{mentorbox}

\section{Move untrusted builds into a safer room}

Today \codepath{LocalFirmwareBuilder} runs Cargo on the host PC. That is suitable for trusted local projects. Cargo can execute build scripts and procedural macros, so an unknown ZIP deserves a stronger boundary.

The \codepath{FirmwareBuilder} trait leaves a doorway for \codepath{DockerFirmwareBuilder}. A careful container would use a pinned toolchain image, a read-only source mount, a small writable output folder, and limits for time, CPU, memory, processes, and logs. It would carry no host secrets.

The coordinator can keep receiving the same log lines and final artifact. Routes and frontend polling do not need to learn how the build room changed.

\section{Give firmware a real signature}

Keep SHA-256 for transfer integrity. Add publisher proof as a separate layer built from supported signing tools and device trust anchors.

The build system would sign an approved image or digest with a protected private key. Firmware metadata would carry the signature, algorithm, and key ID. The device would verify them before activating the slot. Key rotation and recovery need their own tested plan.

Use ESP-IDF Secure Boot and signing guidance when moving toward production. A homemade shared-secret formula is too fragile for the job.

\section{Let the API grow without surprising old projects}

Managed uploads already declare \codepath{application_api_version = 2}. When the application contract changes, add or migrate versions deliberately. Give an old ZIP a clear rejection message and a short conversion example.

Prefer adding fields that old code can ignore. Keep runtime and application API versions pinned together in the generated workspace. If public HTTP payloads need a breaking change, give the new contract a visible version boundary.

\section{Test from the cheap layer toward the board}

Start with unit tests for paths, versions, manifest reading, and status changes. Add integration tests around a temporary SQLite database and a fake builder. Check frontend navigation and failure states in the browser. Then run a known-good project through the real ESP toolchain.

Physical tests come last because they take the most time and reveal a different class of problem. Test a successful update, unplugged network, interrupted download, wrong hash, power loss, setup failure, rollback, and repeated slot swapping.

Write down which layers passed. ``The Rust tests pass'' and ``rollback worked on this board'' are different pieces of evidence.

\begin{checkpoint}
Pick one extension from this chapter. Fill in the ownership sentence at the beginning. Below it, draw the files you would touch, the new failure a learner would see, and the smallest test that proves the change without hardware. Add the physical test on the last line.
\end{checkpoint}



\appendix
\chapter{HTTP API Reference}

This appendix is a desk reference. You do not need to read it from top to bottom before using the project. Return here when you are writing a request, checking a field name, or wondering which service owns an address. The story behind these calls lives in Chapters 3 through 9.

\section{Conventions}

The dashboard uses JSON for ordinary requests and responses, multipart form data for project upload, server-sent events for CSV progress, HTML for profiling reports, and raw octet streams for firmware downloads.

OTA failures share this envelope:

\begin{lstlisting}[language=JSON,numbers=none]
{
  "success": false,
  "error": {
    "code": "TARGET_MISMATCH",
    "message": "Device chip does not match firmware target",
    "details": "optional diagnostic detail",
    "hint": "optional corrective action"
  }
}
\end{lstlisting}

Fields with no value may be omitted. IDs are UUID strings except \codepath{device_id}, which is chosen by configuration or derived from the station MAC.

\section{OTA endpoint index}

{\small
\begin{tabularx}{\textwidth}{>{\raggedright\arraybackslash}p{0.1\textwidth}>{\raggedright\arraybackslash}p{0.52\textwidth}Y}
\toprule
Method & Endpoint & Purpose\\
\midrule
GET & \codepath{/api/ota/status} & Service health and public URL\\
GET & \codepath{/api/ota/projects} & List uploaded projects\\
POST & \codepath{/api/ota/projects} & Upload a project ZIP\\
POST & \codepath{/api/ota/projects/<project-id>/build} & Queue a build\\
GET & \codepath{/api/ota/builds} & List build history\\
GET & \codepath{/api/ota/builds/<build-id>} & Build status and full log\\
GET & \codepath{/api/ota/firmware} & List ready firmware\\
GET & \codepath{/api/ota/firmware/latest?target=esp32s3} & Newest ready target image\\
GET & \codepath{/api/ota/firmware/<firmware-id>} & Firmware details\\
GET & \codepath{/api/ota/firmware/<firmware-id>/download} & Stream application image\\
DELETE & \codepath{/api/ota/firmware/<firmware-id>} & Delete unreferenced image\\
GET & \codepath{/api/ota/devices} & List registered devices\\
POST & \codepath{/api/ota/devices/register} & Create or update a device\\
POST & \codepath{/api/ota/devices/<device-id>/heartbeat} & Refresh device state\\
POST & \codepath{/api/ota/devices/<device-id>/deploy} & Select desired firmware\\
GET & \codepath{/api/ota/devices/<device-id>/update} & Poll desired firmware\\
POST & \codepath{/api/ota/devices/<device-id>/ota-status} & Report OTA progress/result\\
\bottomrule
\end{tabularx}
}

\section{Status}

\textbf{GET \codepath{/api/ota/status}} returns 200.

\begin{lstlisting}[language=JSON,numbers=none]
{
  "service": "CSV_PROFILING OTA Server",
  "language": "rust",
  "status": "online",
  "version": "0.1.0",
  "public_base_url": "http://192.168.1.5:7000",
  "supported_targets": ["esp32s3"]
}
\end{lstlisting}

\section{Projects and builds}

\textbf{POST \codepath{/api/ota/projects}} expects multipart fields \codepath{project}, \codepath{target}, and \codepath{version}. A successful upload returns 201.

\begin{lstlisting}[numbers=none]
curl -F "project=@uploaded-firmware-basic.zip" \
  -F "target=esp32s3" -F "version=1.0.1" \
  http://localhost:7000/api/ota/projects
\end{lstlisting}

\begin{lstlisting}[language=JSON,numbers=none]
{
  "project_id": "55e4...",
  "name": "uploaded-firmware-basic",
  "target": "esp32s3",
  "version": "1.0.1",
  "status": "uploaded",
  "project_type": "managed_application",
  "application_api_version": 2,
  "entrypoint": "src/main.rs"
}
\end{lstlisting}

Standalone uploads omit the API version and entrypoint and report \codepath{project_type} as \codepath{standalone}.

\textbf{GET \codepath{/api/ota/projects}} returns an array of project records. Internal filesystem paths are skipped during serialization.

\textbf{POST \codepath{/api/ota/projects/<project-id>/build}} returns 202:

\begin{lstlisting}[language=JSON,numbers=none]
{"build_id":"a8c4...","status":"building"}
\end{lstlisting}

The durable row begins queued and is moved to building by the coordinator task. Only one queued or active build is allowed for a project.

\textbf{GET \codepath{/api/ota/builds}} returns all builds, newest first. \textbf{GET \codepath{/api/ota/builds/<build-id>}} returns one response shaped like:

\begin{lstlisting}[language=JSON,numbers=none]
{
  "build_id": "a8c4...",
  "project_id": "55e4...",
  "project_name": "uploaded-firmware-basic",
  "target": "esp32s3",
  "version": "1.0.1",
  "status": "success",
  "logs": ["Build queued", "Build started", "Tooling: cargo ..."],
  "exit_code": 0,
  "error": null,
  "started_at": "2026-09-07T10:00:00Z",
  "finished_at": "2026-09-07T10:02:00Z",
  "created_at": "2026-09-07T09:59:59Z"
}
\end{lstlisting}

Failed responses may include \codepath{hint}.

\section{Firmware}

List, detail, and latest routes return the same firmware representation:

\begin{lstlisting}[language=JSON,numbers=none]
{
  "firmware_id": "c53b...",
  "project_id": "55e4...",
  "build_id": "a8c4...",
  "project": "uploaded-firmware-basic",
  "target": "esp32s3",
  "version": "1.0.1",
  "filename": "firmware-uploaded-firmware-basic-v1.0.1-a8c4.bin",
  "size": 945332,
  "sha256": "f2534c...",
  "status": "ready",
  "created_at": "2026-09-07T10:02:00Z",
  "download_url": "http://192.168.1.5:7000/api/ota/firmware/c53b.../download"
}
\end{lstlisting}

\textbf{GET \codepath{/download}} returns \codepath{application/octet-stream}, content length, attachment disposition, and \codepath{X-Firmware-SHA256}. The route reconstructs the path from the configured build directory and stored filename, then canonicalizes it before opening.

\textbf{DELETE \codepath{/api/ota/firmware/<firmware-id>}} returns 204. It returns 409 when a device or deployment references the image.

\section{Devices and deployments}

\textbf{POST \codepath{/api/ota/devices/register}} returns 201 and the current device record.

\begin{lstlisting}[language=JSON,numbers=none]
{
  "device_id": "ESP32-S3-001",
  "name": "Sensor Node 1",
  "chip": "esp32s3",
  "current_version": "1.0.0",
  "ip_address": "192.168.1.24"
}
\end{lstlisting}

\textbf{POST \codepath{/api/ota/devices/ESP32-S3-001/heartbeat}} returns 200 and the updated device.

\begin{lstlisting}[language=JSON,numbers=none]
{
  "current_version": "1.0.0",
  "status": "online",
  "ip_address": "192.168.1.24"
}
\end{lstlisting}

Allowed heartbeat statuses are online, offline, updating, rebooting, updated, failed, and unknown.

\textbf{GET \codepath{/api/ota/devices}} returns records with these public fields:

\begin{lstlisting}[language=JSON,numbers=none]
{
  "device_id": "ESP32-S3-001",
  "name": "Sensor Node 1",
  "chip": "esp32s3",
  "ip_address": "192.168.1.24",
  "current_version": "1.0.0",
  "desired_version": "1.0.1",
  "desired_firmware_id": "c53b...",
  "status": "online",
  "ota_status": "pending",
  "ota_progress": 0,
  "ota_error": null,
  "last_seen": "2026-09-07T10:03:00Z"
}
\end{lstlisting}

\textbf{POST \codepath{/api/ota/devices/<device-id>/deploy}} expects a firmware ID and returns 202.

\begin{lstlisting}[language=JSON,numbers=none]
{"firmware_id":"c53b..."}
\end{lstlisting}

\begin{lstlisting}[language=JSON,numbers=none]
{
  "deployment_id": "d911...",
  "device_id": "ESP32-S3-001",
  "firmware_id": "c53b...",
  "desired_version": "1.0.1",
  "status": "pending"
}
\end{lstlisting}

\textbf{GET \codepath{/api/ota/devices/<device-id>/update}} returns either \codepath{update_available: false} or the manifest shown in Chapter 9.

\textbf{POST \codepath{/api/ota/devices/<device-id>/ota-status}} accepts:

\begin{lstlisting}[language=JSON,numbers=none]
{
  "status": "downloading",
  "progress": 35,
  "current_version": null,
  "error": null
}
\end{lstlisting}

Allowed OTA states are idle, pending, downloading, verifying, installing, rebooting, success, and failed. Progress must be between 0 and 100. Failed requires an error. Success defaults to 100 and clears desired firmware after updating the current version.

\section{CSV endpoint index}

These routes belong to Flask on port 5000.

\begin{longtable}{p{0.13\textwidth}p{0.38\textwidth}p{0.39\textwidth}}
\toprule
Method & Endpoint & Purpose\\
\midrule
\endhead
GET & \codepath{/} & Serve the dashboard HTML\\
POST & \codepath{/upload} & Accept multipart CSV field \codepath{file} and return \codepath{job_id}\\
GET & \codepath{/progress/<job-id>} & Stream profiling progress with SSE\\
GET & \codepath{/report/<job-id>} & Serve generated report HTML\\
GET & \codepath{/sessions} & List retained report sessions\\
GET & \codepath{/session/<job-id>} & Check and restore one session\\
POST & \codepath{/clean/<job-id>} & Apply cleaning controls and refresh metadata\\
POST & \codepath{/playground/<job-id>} & Run preview, analysis, chart, or explicit mutation\\
GET & \codepath{/download-cleaned/<job-id>} & Download cleaned CSV\\
DELETE & \codepath{/reset/<job-id>} & Delete one CSV session\\
DELETE & \codepath{/reset-all} & Delete all CSV sessions\\
\bottomrule
\end{longtable}

Playground request fields include \codepath{operation}, \codepath{mutate}, \codepath{limit}, \codepath{column}, \codepath{column_2}, \codepath{value}, \codepath{operator}, \codepath{indexes}, and \codepath{chart_type}. The endpoint clamps limit to 1 through 500.

\chapter{Configuration, Commands, and Glossary}

Keep this appendix near your terminal. It gathers names and commands that are easy to forget after you understand the larger idea. Copy a command only after reading the short note around it, especially when hardware is connected.

\section{Environment variables}

\begin{longtable}{>{\ttfamily\small}p{0.29\textwidth}p{0.22\textwidth}p{0.39\textwidth}}
\toprule
Variable & Default/example & Meaning\\
\midrule
\endhead
OTA\_BIND\_ADDRESS & \codepath{0.0.0.0} & Listener interface; LAN access needs a non-loopback binding.\\
OTA\_PORT & \codepath{7000} & Axum HTTP port.\\
OTA\_PUBLIC\_BASE\_URL & detected or explicit & Base placed in device-facing firmware URLs. Use the PC LAN address.\\
OTA\_PROJECT\_DIR & \codepath{firmware-projects} & Safe root for extracted uploads.\\
OTA\_BUILD\_DIR & \codepath{firmware-builds} & Safe root for generated application images.\\
OTA\_DATA\_DIR & \codepath{data} & Directory containing \codepath{ota.db}.\\
OTA\_RUNTIME\_DIR & \codepath{examples/esp32-rust-ota-client} & Stable runtime template used for managed builds.\\
OTA\_MAX\_UPLOAD\_MB & \codepath{50} & Encoded project field limit in MiB.\\
OTA\_MAX\_EXTRACTED\_MB & upload limit times 10 & Maximum total declared extracted content.\\
OTA\_ALLOWED\_ORIGINS & localhost and 127.0.0.1 on port 5000 & Comma-separated browser origin allow-list.\\
OTA\_RUST\_TOOLCHAIN & \codepath{esp} & Rustup toolchain applied to builder commands; empty uses the process default.\\
WIFI\_SSID & none & Embedded Wi-Fi network name used at managed build time.\\
WIFI\_PASS & none & Embedded Wi-Fi password. Keep it in local \codepath{.env}.\\
DEVICE\_NAME & managed-device label & Friendly device name embedded into managed firmware.\\
DEVICE\_ID & MAC-derived when absent & Optional stable identity override.\\
RUST\_LOG & info defaults & Axum and tower-http tracing filter.\\
\bottomrule
\end{longtable}

\section{Command sheet}

\subsection{Start services}

\begin{lstlisting}[numbers=none]
# Windows
.\run-dev.bat

# Linux/macOS
./run-dev.sh
\end{lstlisting}

\subsection{Run OTA server checks}

\begin{lstlisting}[numbers=none]
cargo test --manifest-path ota-server/Cargo.toml
cargo clippy --manifest-path ota-server/Cargo.toml --all-targets
curl http://localhost:7000/api/ota/status
\end{lstlisting}

\subsection{Check embedded tools}

\begin{lstlisting}[numbers=none]
cargo --version
rustc +esp --version
espflash --version
rustc +esp --print target-list
\end{lstlisting}

\subsection{Build the stable runtime without flashing}

\begin{lstlisting}[numbers=none]
cd examples/esp32-rust-ota-client
cargo +esp build --release --target xtensa-esp32s3-espidf
\end{lstlisting}

\subsection{Build and flash the stable runtime}

\begin{lstlisting}[numbers=none]
cargo +esp run --release --target xtensa-esp32s3-espidf
\end{lstlisting}

The runner settings determine the exact \codepath{espflash} arguments and partition-table use. Inspect \codepath{.cargo/config.toml} in the embedded project before physical flashing.

\section{Glossary}

\begin{description}
  \item[Application image] A \codepath{.bin} containing one ESP-IDF application, suitable for an OTA app partition.
  \item[Axum] The Rust HTTP framework used by the OTA service.
  \item[Build ID] UUID identifying one compilation attempt and its durable logs.
  \item[Cargo target] The architecture and platform a Rust crate is compiled for. This project begins with \codepath{xtensa-esp32s3-espidf}.
  \item[CORS] Browser policy controlling which page origins may call the OTA API.
  \item[Deployment] A stored relationship saying one device should install one firmware record.
  \item[ELF] Linked executable format consumed by \codepath{espflash save-image}.
  \item[ESP-IDF] Espressif's official development framework and system services used beneath Rust bindings.
  \item[Firmware ID] UUID identifying one ready application image and metadata record.
  \item[Flask] The Python HTTP framework serving the dashboard and CSV API.
  \item[Heartbeat] Periodic device request refreshing version, status, IP address, and last-seen time.
  \item[Inactive slot] OTA application partition that is not currently executing and can safely receive a new image.
  \item[Managed application] Uploaded function-based Rust package composed with the stable runtime by the server.
  \item[NVS] ESP-IDF non-volatile storage, commonly used for settings and Wi-Fi configuration.
  \item[OTA] Over-the-air firmware update performed through a network connection.
  \item[Pull model] The device polls for desired firmware and initiates its own download.
  \item[Rollback] Bootloader-assisted return to a previously valid application slot when a new image fails startup validation.
  \item[SHA-256] Hash used here to confirm downloaded bytes match server metadata.
  \item[SSE] Server-sent events, the one-way HTTP stream used for CSV profiling progress.
  \item[Stable runtime] Initially flashed platform code that owns Wi-Fi, OTA, rollback, and application lifecycle.
  \item[Standalone firmware] Uploaded Cargo project built without managed-runtime composition; it owns its complete device behavior.
  \item[WAL] SQLite write-ahead logging mode, used to improve local concurrent access behavior.
  \item[ZIP slip] Archive path traversal that attempts to extract outside the intended root.
\end{description}

\section{Current implementation limits}

The first version supports ESP32-S3, local SQLite storage, polling for build logs, plain HTTP on a trusted LAN, local trusted builds, and SHA-256 integrity checks. Device serial/network terminal streaming, isolated Docker builds, multi-user authentication, publisher signing, and additional chip targets are extension work described in Chapter 12.

These limits are design boundaries, not hidden promises. Keeping them visible helps a beginner distinguish what can be used today from what still needs engineering and hardware validation.


\backmatter
\chapter*{Closing Note}
\addcontentsline{toc}{chapter}{Closing Note}
You now have a map from browser click to Pandas report, and from uploaded Rust source to an application image written into an inactive ESP32 OTA slot. That map matters more than memorizing a command. When something fails, locate the boundary first: browser to Flask, browser to Axum, builder to toolchain, PC to LAN, or bootloader to application. The logs on that boundary usually tell you what to fix next.

Keep the initial runtime recoverable over USB, test rollback on real hardware, and treat uploaded build inputs as trusted code until the builder runs inside a sandbox. With those habits, this project can grow without losing the clean separation that makes it understandable.

\end{document}
